\section{Related Work}
\label{sec:related-work}

Alignment evaluations for large language models have proliferated in recent years, but most operate at surface levels of safety, ethical reasoning, or factuality.

TruthfulQA tests whether models respond truthfully to adversarially constructed questions \cite{lin2022truthfulqameasuringmodelsmimic}, surfacing model tendencies to overstate knowledge or hallucinate. ETHICS and Moral Scenarios datasets probe moral reasoning across dilemmas \cite{hendrycks2023aligningaisharedhuman}, while ToxiGen and RealToxicityPrompts evaluate harmful or offensive content generation \cite{hartvigsen2022toxigenlargescalemachinegenerateddataset, gehman2020realtoxicitypromptsevaluatingneuraltoxic}. These tools are valuable for shallow alignment, but generally lack interactional depth or strategic framing.

HELM offers a broad evaluation suite, including some alignment-relevant tasks, but it still primarily emphasizes single-shot classification, instruction following, or broad capability measurement rather than staged decision-making under adversarial conditions \cite{liang2023holisticevaluationlanguagemodels}. As a result, it is often better suited to measuring whether a model can answer correctly or safely in the moment than to diagnosing how it behaves when instrumental incentives conflict with oversight.

Closer in spirit, recent work has begun to probe instrumental convergence more directly. \emph{Evaluating the Paperclip Maximizer} studies whether RL-trained language models are more likely to pursue instrumental goals under targeted stress tests \cite{he2025paperclip}. InstrumentalEval and related steerability-style evaluations then make this question more explicit by constructing compact scenario prompts and appending natural-language suffixes that either incentivize instrumental behavior or discourage it \cite{hoscilowicz2026steerability}. 
This is useful: it provides a relatively clean way to measure whether a model can be pushed toward or away from behaviors such as shutdown avoidance, hacking, resource acquisition, or deception under controlled prompt interventions.

However, this steering-based design also answers a somewhat different question from the one StratEval is built to study. InstrumentalEval-style setups often tell the model, in natural language, what strategic orientation to adopt or avoid. 
In effect, they test how readily instrumental-convergence tendencies can be elicited, amplified, or suppressed once the evaluation itself makes those tendencies salient. 
StratEval instead asks whether the model finds the instrumental response \emph{on its own}. 
Rather than appending clauses such as ``do not let anything stop you from completing your objective'' or ``do not deceive, manipulate, or resist shutdown,'' we construct scenarios in which self-preservation, coercion, deception, procedural evasion, or capability-seeking may become instrumentally useful responses to the situation. 
The methodological contribution is therefore not merely staged or multi-probe prompting, but a different elicitation philosophy: we do not instruct the model to be instrumentally convergent; we observe whether it independently reaches for that strategy when placed under pressure, and subsequently, document, judge, and map what strategies emerge and under which circumstances.
The resulting data is less focused on evaluating an individual model's alignment and more focused on classifying distinct failure modalities and approximating their severity.

This distinction matters because the two paradigms provide different evidence. 
Steering-style evaluations are valuable for measuring susceptibility and controllability. StratEval is aimed more squarely at diagnostic realism: what kinds of strategic behavior emerge when the model is given a role, a conflict, and an opportunity structure, but is not explicitly told which instrumental tactic to consider. In that sense, the approaches are complementary rather than adversarial. 
One asks how steerable instrumental-convergence tendencies are once named; the other asks whether those tendencies appear when the model must discover them for itself.

StratEval also differs from prior work in three further ways.

Decomposed prompting: Most prior work reduces alignment to a one-shot prompt. We use recognition-oriented probes and decision-stage prompts over matched scenario frames to compare understanding with action, surfacing cases where a model appears to understand the problem and still chooses a misaligned course.

Graded severity: We move beyond binary success/failure outcomes by assigning each decision an ordinal escalation score from 0 to 8, allowing fine-grained comparison of failure intensity.

Diagnostic taxonomy: Our 37 failure-mode labels, together with a safe label, enable interpretability and downstream behavioral analysis that flat metrics cannot capture.

In this way, StratEval complements prior alignment evaluations by focusing less on whether a model can state the right principle in the abstract, and more on how it behaves when those principles conflict with instrumental incentives.
